\documentclass[xcolor=pdftex,svgnames,table]{beamer}
\usepackage{latexsym,amssymb,amsmath,stmaryrd}
\usepackage[utf8]{inputenc}
\usepackage[T1]{fontenc}
\usepackage[francais]{babel}
\usepackage{pgf}
\usepackage{tikz}
\usetikzlibrary{arrows,patterns,plotmarks,shapes,snakes,er,3d,automata,%
backgrounds,topaths,trees,petri,mindmap}
\usepackage{pgfbaseimage}
\usepackage{ifpdf}
\ifpdf
\usepackage{graphicx}
\else
%\usepackage[dvips]{graphicx}
\usepackage{pstricks,pst-tree,pst-node}
\newpsobject{showgrid}{psgrid}{subgriddiv=1,griddots=10,gridlabels=6pt}
\fi
\usepackage{moreverb}
\usepackage[normalem]{ulem}
\usepackage{version}
\usepackage{url}
\usepackage{multicol}
\usepackage{listings}
\lstset{
  language=[ANSI]C, 
%  gobble=2, 
%  escapeinside="", 
  basicstyle=\ttfamily, 
%  directivestyle=\color{Fuchsia},
  identifierstyle = \color{DarkOrange},
  keywordstyle=\color{DarkGreen}, 
  commentstyle=\color{red}, 
%  numbers=left, 
%  numbersep=5pt, 
%  numberstyle=\scriptsize, 
%  moredelim=[is][\color{gray}\itshape]{/*}{*/}, 
%  moredelim=[is][\alert]{/+}{+/}, 
%  morecomment=[is]{/=}{=/}, 
%  morecomment=[in]{comment=}{,} 
%  fancyvrb=true 
} 

\newcommand{\binaire}[1]{\ensuremath{\underline{#1}}}
\newcommand{\C}[1]{\texttt{#1}}
\newcommand{\bbbn}{\ensuremath{\mathbb{N}}}
\input{communs/cc_beamer}
\usetheme{classic}
\newcommand{\nowrite}{\put(10,-4){\includegraphics[scale=.05]{creative_commons/nopencil}}}
\newcommand{\youwrite}{\put(10,-4){\includegraphics[scale=.05]{creative_commons/pencil}}}
\newcommand{\aemporter}{\put(10,-6){\includegraphics[scale=.05]{creative_commons/szymonraj_Shopping_bag}}}


%%% Titre -- cours 8
\title{Bases de programmation -- Cours 4.\\ Fonctions et procédures}
\author{Pierre Boudes}
\date{\today}

\begin{document}


%% Page de titre et licence CC.
\begin{frame}
	\titlepage
	\vfill
	\begin{center}
		\CcGroupByNcSa{0.83}{0.95ex}\\[2.5ex]
		{\tiny\CcNote{\CcLongnameByNcSa}}
		\vspace*{-2.5ex}
	\end{center}
\end{frame}


%%%%%%%%%%%%%%%%%%%%

\begin{frame}
  \frametitle{Cours en plusieurs parties sur les fonctions \nowrite}

\begin{enumerate}
    \item Principe de base, intérêt et analogie mathématique
      \item \emph{Déclarer, appeler, définir} ou \emph{appeler, déclarer, définir}
    \item Fonctions sans entrée ou sans sortie, effets de bord
    \item Pile d'appel, fonctions récursives
    \item Cours suivant : autres types que les entiers pour les paramètres et la sortie
  \end{enumerate}
\end{frame}

\section[Plan]{}
\frame[label=plan]{\frametitle{Plan de la séance\nowrite}\tableofcontents}


\section[Introduction]{Définition et intérêt des fonctions en informatique}
\begin{frame}
  \frametitle{Définition et intérêt des fonctions en informatique}
  \begin{quote}
    En informatique, une fonction est une portion de code représentant
    un sous programme, qui effectue une tâche ou un calcul
    relativement indépendant du reste du programme. (wikipédia)
  \end{quote}
\pause
\begin{quote}
  Les fonctions c'est de la sous-traitance en interne.
\end{quote}
\pause
Intérêt des fonctions :
  \begin{itemize}
  \item \structure{factorisation} : éviter de recopier du code.\pause
    \item \structure{réutilisation} : copier-coller dans d'autres
      programmes, utiliser des bibliothèques (\C{rand()}, etc.);\pause
    \item \structure{lisibilité} : faciliter la relecture du code et cacher les détails;\pause
    \item \structure{structuration} : découper en sous-problèmes de
      tailles pensables;  distribuer la programmation à différente personnes, à
      différentes étapes de la réalisation d'un projet.\pause
  \end{itemize}
\end{frame}

\subsection{Analogie mathématique}
\begin{frame}
  \frametitle{Analogie mathématique \nowrite}
 \begin{itemize}
\item De même que le terme \emph{variable} a un sens
    différent en informatique et en mathématiques, le terme
    \emph{fonction} recouvre des réalités différentes.

\item Pour les fonctions, les deux notions sont historiquement très
    proches. C'est bien pour écrire des fonctions mathématiques que
    les fonctions informatiques ont été introduites.

\item Mais qu'est-ce qu'une fonction mathématique ?
 \end{itemize}
\end{frame}

\subsection{La notion mathématique de fonction}
\begin{frame}
  \frametitle[En maths]{La notion mathématique de
    fonction}
 Une fonction mathématique est :\pause
  \begin{itemize}
  \item Une courbe ?
  \item Une expression avec inconnues ($x^2 + 1$), une formule ?
  \item Un terme dans une algèbre de fonctions ($f\circ g$, $f'$)
    ?\pause
  \item Une relation entre deux ensembles, entrées et sorties, telle
    que chaque entrée est en relation avec au plus une sortie
    (toujours la même). Le graphe de la fonction.
  \end{itemize}
\pause
\begin{block}{À retenir}
  En mathématiques, la manière dont la fonction est calculée ne fait
  pas partie de l'objet défini. Une fonction est une boîte noire.
\end{block}

\end{frame}

\section[En C]{Les fonctions en langage C}
\begin{frame}
\frametitle{Déclarer, appeler, définir.}
\pause
\begin{itemize}
\item \structure{Déclarer} : comme les variables, les fonctions
  doivent être déclarées avant usage pour fixer \alert{le nombre} et le type
  des arguments et de la sortie.\pause
\item \structure{Appeler} : utiliser une fonction, faire appel à son
  résultat en fixant les valeurs des arguments (paramètres effectifs).\pause
\item \structure{Définir} : décrire le corps de la fonction, c'est à
  dire la suite d'instructions qui constitue son calcul (sur des
  paramètres formels).
\end{itemize}\pause
\begin{description}
\item[Ordre de la lecture (compilation) :] déclaration puis appels et définition
\item[Ordre de l'écriture  (création) :] appels puis déclaration et définition
\end{description}


\end{frame}

\subsection{Appeler des fonctions}
\begin{frame}[fragile]
  \frametitle{Utiliser des fonctions : appeler}

Supposons que l'on dispose d'une fonction \C{maximum}, qui lorsqu'on
lui donne deux entiers nous renvoie le plus grand des deux.
\pause
Nous pouvons l'utiliser de différentes façons :
\begin{lstlisting}[escapechar={\%},basicstyle=\ttfamily\small]
int main ()
{
  int x = 3;
  int y = 4;
  int z;

  z = %\only<2->{\colorbox{yellow}{maximum}}%(x, y);%\pause%
  z = 2 + maximum(3 * 2, z);%\pause%
  x = maximum(maximum(3, y) + 1, z - 1);
\end{lstlisting}
\pause
Chaque entrée prend une valeur en fonction de l'expression passée en
argument. La fonction prend alors sa valeur de sortie.
%\pause
%\structure{Remarque :} \C{printf} et \C{scanf} sont des fonctions (un peu spéciales).
\end{frame}

\subsection{Créer des fonctions : déclarer, définir}
\begin{frame}[fragile]
  \frametitle{Créer des fonctions (1)  déclarer}
Déclarer une fonction c'est donner son prototype (ou signature), c'est
à dire son nom, le type de ses arguments, le type de sa sortie.\pause
\begin{lstlisting}[basicstyle=\ttfamily\scriptsize]
type_sortie nom_fonction(type1 param1, ..., typen paramn);
\end{lstlisting}
\pause
\begin{block}{Remarque}
  Seule la déclaration des types est importante pour
  l'analyse sémantique. Cependant, le choix de noms de paramètres
  explicites fournit une documentation minimale.
\end{block}


\pause
\begin{block}{Exemple}

\begin{lstlisting}[basicstyle=\ttfamily\small]
int puissance(int, int);
int puissance(int base, int exposant);
\end{lstlisting}
\end{block}
\pause
\textbf{La déclaration doit se trouver avant tout appel.}
\end{frame}
\begin{frame}[fragile]
  \frametitle{Créer des fonctions (2)  définir}
\pause
Définir une fonction c'est donner l'ensemble des instructions qui
permettent, à partir des paramètres d'entrée, de calculer la valeur de
la fonction, dite valeur de retour, valeur renvoyée.\pause
\begin{lstlisting}[basicstyle=\ttfamily\scriptsize]
type_sortie nom_fonction(type1 param1, ..., typen paramn)
{
  /* declaration et initialisation variables */

  /* instructions */
}
\end{lstlisting}\pause
Les noms des paramètres
formels sont \C{param1}, \ldots, \C{paramn}, ils sont utilisés dans le
corps du calcul.
\pause

Dès que la valeur de la
fonction est calculée, on utilise l'instruction \C{return} qui marque
la fin du calcul et donne sa valeur à la fonction. \pause
\begin{lstlisting}[basicstyle=\ttfamily\small]
  return expression_resultat;
\end{lstlisting}

La définition est obligatoire pour que l'édition de liens réussise et que l'exécutable soit créé.
\end{frame}
\begin{frame}
    \frametitle{Variables locales, portée, durée}

    Le corps de la fonction peut déclarer des variables additionnelles,
    qui sont locales à la fonction (portée) et se voient allouer un
    espace mémoire pour chaque appel de la fonction (durée).

    \begin{block}{Remarque} \C{main} est une fonction (appelée au lancement du
      programme).
    \end{block}
\end{frame}


\begin{frame}
  \frametitle{Résumé}
 Utilisation des fonctions :
  \begin{itemize}
    \item \structure{déclaration} (types des paramètres et de la valeur de retour)
    \item \structure{définition}  (code, paramètres formels)
    \item \structure{appel} (paramètres effectifs, espace mémoire)
  \end{itemize}
\pause
\begin{block}{Convention de nommage}
  Il est pratique de faire commencer le nom de chaque fonction par un
  verbe à l'infinitif \C{convertir\_\ldots}, \C{tester\_\ldots},
  etc. ou éventuellement par un verbe conjugué : par exemple \C{est\_}
  pour des fonctions à valeur de test (\C{est\_majeur},
  \C{est\_premier}, \ldots...)

Exceptions : la fonction est connue sous un autre nom
  (\C{factorielle(n)}, \C{racine(x)}, \C{moyenne(...)})
\end{block}
\end{frame}




\subsection[Procédures]{Fonctions sans valeurs de retour (void)}
\begin{frame}[fragile]
  \frametitle{Fonctions sans valeurs de retour (void)}

On parle plutôt de procédure ou de routine car l'analogie avec les
fonctions mathématiques est perdue.

 \begin{block}{Déclarer}
    \begin{lstlisting}[basicstyle=\ttfamily\small]
void afficher_valeurs(int x, int y);
     \end{lstlisting}
  \end{block}

  \begin{block}{Appeler}
  \begin{lstlisting}[basicstyle=\ttfamily\small]
afficher_valeurs(5, 3);
   \end{lstlisting}
  \end{block}

  \begin{block}{Définir}
Comme d'habitude mais pas de \C{return} (ou \C{return sans argument}).
\end{block}
\end{frame}

%\subsection[Procédures]{Fonctions sans valeurs de retour (void)}
\begin{frame}[fragile]
  \frametitle{Fonctions sans arguments}

 \begin{block}{Déclarer}
    \begin{lstlisting}[basicstyle=\ttfamily\small]
int nombre_aleatoire();
int saisie_utilisateur();
     \end{lstlisting}
  \end{block}

  \begin{block}{Appeler}
  \begin{lstlisting}[basicstyle=\ttfamily\small]
  int n;
  int secret;

  secret = nombre_aleatoire();

  n = saisie_utilisateur();
   \end{lstlisting}
  \end{block}

  \begin{block}{Définir}
Comme d'habitude.
\end{block}
\end{frame}

\subsection[math.h]{Utiliser les fonctions d'une bibliothèque}

\begin{frame}[fragile]
  \frametitle{Utiliser les fonctions d'une bibliothèque (math.h)}
Utilisation de la bibliothèque math.h
\begin{verbatim}
$ man math
\end{verbatim}
  \begin{block}{Déclarer}
    \begin{lstlisting}[basicstyle=\ttfamily\small]
#include <math.h>
     \end{lstlisting}
  \end{block}

  \begin{block}{Appeler}
  \begin{lstlisting}[basicstyle=\ttfamily\small]
  double x; /* nombre a virgule (cours 5) */

  x = log(3.5); /* la virgule se note avec un point */
    \end{lstlisting}
  \end{block}

  \begin{block}{Définir}
\begin{verbatim}
$ gcc -lm -Wall prog.c -o prog.exe
\end{verbatim}
 \end{block}
\end{frame}

%\subsection{Rappel sur les fonctions en C}
\section{Traces et pile d'appel}

\begin{frame}
  \frametitle{Traces et pile d'appel\nowrite}
 Utilisation des fonctions :
  \begin{itemize}
    \item \structure{déclaration} \textbf{types} (\C{int}, \C{char},
      \C{double}, \C{void}) des paramètres et de la valeur de retour;
    \item \structure{définition}  code, paramètres formels (chacun est
      une variable locale);
    \item \structure{appel} paramètres effectifs (chaque expression
      donnant sa valeur au paramètre formel correspondant), \textbf{espace mémoire}.
  \end{itemize}

Nous allons voir de façon plus précise cette question
d'espace mémoire, avec la \emph{pile d'appel}…

\end{frame}



\subsection{Traces : flot de contrôle et données}
\begin{frame}
  \frametitle{Traces}

Pour étudier l'exécution de nos programmes, nous simulons leur
exécution à la main. Comment tenir compte de chaque appel de
chaque fonction que l'on a défini ?

{
\rowcolors{2}{LightSkyBlue}{LightSteelBlue} %\arrayrulecolor{green!75!gray}
\small
      \setlength{\unitlength}{\tabcolsep}
          \begin{tabular}[t]{|c|c|l|}
          \multicolumn{3}{l}{\C{main()}}\\ \hline
          ligne & n & Affichage \\ \hline
          initialisation  & 9 & \\ \hline
          16 & & \\ \hline
          \rowcolor{LightYellow}
          & & \multicolumn{1}{>{\columncolor{LightYellow}}r|}{
            \put(1,0){\noindent
              \begin{tabular}[t]{|c|c|c|l|}
          \multicolumn{4}{l}{\C{est\_premier(9)}}\\ \hline
          ligne & n & i & Affichage \\  \hline
          initialisation  & 9 & ? & \\ \hline
         34 &  & 2 & \\ \hline
         40  &  & 3 & \\ \hline
         38  &\multicolumn{3}{|>{\columncolor{MediumSpringGreen}}l|}{renvoie false}\\ \hline
        \end{tabular}
              }
          }\\ \hline
          22 & & 9 n'est pas premier\\ \hline
          26 &\multicolumn{2}{|>{\columncolor{LightSteelBlue}}l|}{renvoie EXIT\_SUCCESS}\\ \hline
          \end{tabular}
}

On ne \emph{trace} pas les fonctions externes, comme \C{printf}.
\end{frame}


\begin{frame}[fragile,label=current1]
  \frametitle{Focus sur le passage de valeurs}
  Notez que les fonctions communiquent \textbf{des valeurs}, pas des noms de
  variables.
\begin{columns}
\scriptsize
  \begin{column}[t]{3.7cm}
 \begin{lstlisting}[numbers=left,basicstyle=\ttfamily\scriptsize]
void permute_valeurs(int a,int b);

int main() {
    int x = 1;
    int y = 2;
    permute_valeurs(x,y);
    printf("x = %d et y = %d\n",x,y);
    return EXIT_SUCCESS;
}

void permute_valeurs(int a,int b) {
    int aux;
    aux = a;
    a = b;
    b = aux;
}\end{lstlisting}
\vspace{.4cm}
 \end{column}
\begin{column}[t]{4cm}
\pause
\rowcolors{2}{LightSkyBlue}{LightSteelBlue} %\arrayrulecolor{green!75!gray}
        \setlength{\unitlength}{\tabcolsep}
      \hspace{-5.5cm} \begin{tabular}[t]{|r|c|c|l|}
          \multicolumn{4}{l}{\C{main()}}\\ \hline
          ligne & x & y & Affichage \\ \hline
          initialisation  & 1 & 2 & \\ \hline
          6 &   & &\pause \\ \hline
          \multicolumn{4}{|r|}{
            \put(1,4){\noindent
              \begin{tabular}[t]{|r|c|c|c|l|}
                \multicolumn{5}{|l}{\C{permute\_valeurs(1, 2)}}\\ \hline
                ligne & a & b & aux & Aff. \\ \hline
                initialisation  & 1 & 2 & ? & \\ \hline
                13 &  &  & 1 & \\ \hline
                14 & 2 &  &  & \\ \hline
                15 &  &  1  &  & \\ \hline
                16 &\multicolumn{4}{|>{\columncolor{MediumSpringGreen}}l|}{ne renvoie rien}\\ \hline
              \end{tabular}
            }} \\ \hline
          7  &  &  & \C{x = \pause1 et y = 2} \\ \hline
          8 &\multicolumn{3}{|l|}{SORTIE AVEC SUCCÈS}\\ \hline
        \end{tabular}
  \end{column}
\end{columns}
\end{frame}
\begin{frame}[fragile,label=current]
  \frametitle{Un exemple avec plus d'appels de fonctions}
\vspace{-0.5cm}
\begin{columns}
\scriptsize
  \begin{column}[t]{3.7cm}
 \begin{lstlisting}[numbers=left,basicstyle=\ttfamily\scriptsize]
double maximum(double x, double y);
double valeur_absolue(double x);
int main() {
    double a = -4.2;
    double b = 1.3;
    a = valeur_absolue(a) + valeur_absolue(b);
    printf("%g\n", a);
    return EXIT_SUCCESS;
}
double maximum(double x, double y) {
    if (x > y) {
        return x;
    }
    return y;
}
double valeur_absolue(double x) {
    return maximum(x, -x);
}
\end{lstlisting}
\vspace{0.5cm}
 \end{column}
\begin{column}[t]{4cm}
\pause~
\only<-6>{%
\rowcolors{2}{LightSkyBlue}{LightSteelBlue}
% \arrayrulecolor{green!75!gray}
        \setlength{\unitlength}{\tabcolsep}
      \hspace{-5cm} \begin{tabular}[t]{|r|c|c|l|}
          \multicolumn{4}{l}{\C{main()}}\\ \hline
          ligne & a & b & Aff. \\ \hline
          initialisation  & -4.2 & 1.3 & \\ \hline
          6\pause & & &  \multicolumn{1}{ >{\columncolor{LightSteelBlue}}r|
                        }{\put(2,0){\C{valeur\_absolue(-4.2)}}} \\ \hline
          \multicolumn{4}{|r|}{\only<7>{\cellcolor{LightGreen}}
            \put(1,1){\noindent
              \begin{tabular}[t]{|r|c|}\hline
                ligne  & x    \\ \hline
                init.  & -4.2 \\ \hline
                17 \pause & \multicolumn{1}{ >{\columncolor{LightSteelBlue}}r|
                }{\put(2,0){\C{maximum(-4.2, 4.2)}}}\\ \hline
                \multicolumn{2}{r|}{\only<7>{\cellcolor{LightGoldenrod}}
                  \put(1,1){\noindent
                    \begin{tabular}[t]{|r|c|c|}\hline
                      ligne & x & y \\ \hline
                      init.  & -4.2 & 4.2 \\ \hline
                      14\pause &\multicolumn{2}{|>{\columncolor{MediumSpringGreen}}l|}{renvoie 4.2}\\ \hline
                    \end{tabular}
                  }} \\ \hline
                17 \pause &\multicolumn{1}{|>{\columncolor{MediumSpringGreen}}l|}{renvoie 4.2}\\ \hline
              \end{tabular}
            }} \\ \hline
          \rowcolor{LightSkyBlue} 6 & & &\multicolumn{1}{ >{\columncolor{LightSkyBlue}}r| }{\put(2,0){\C{valeur\_absolue(1.3)}}}\\ \hline
          \multicolumn{4}{|r|}{\only<7>{\cellcolor{LightGreen}}
            \put(1,1){\noindent
              \begin{tabular}[t]{|r|c|}\hline
                ligne  & x    \\ \hline
                init.  & 1.3 \\ \hline
                17 & \multicolumn{1}{ >{\columncolor{LightSteelBlue}}r|
                }{\put(2,0){\C{maximum(1.3, -1.3)}}}\\ \hline
                \multicolumn{2}{|r|}{\only<7>{\cellcolor{LightGoldenrod}}
                  \put(1,1){\noindent
                    \begin{tabular}[t]{|r|c|c|}\hline
                      ligne & x & y \\ \hline
                      init.  & 1.3 & -1.3 \\ \hline
                      12 &\multicolumn{2}{|>{\columncolor{MediumSpringGreen}}l|}{renvoie 1.3}\\ \hline
                    \end{tabular}
                  }} \\ \hline
                17 &\multicolumn{1}{|>{\columncolor{MediumSpringGreen}}l|}{renvoie 1.3}\\ \hline
              \end{tabular}
            }} \\ \hline
          6  & 5.5 &  &  \\ \hline
          7  &  &  & \C{5.5} \\ \hline
          8 &\multicolumn{3}{|l|}{SUCCÈS}\\ \hline
        \end{tabular}
}%fin only
\only<7->{%
\rowcolors{2}{LightGreen}{LightGreen}
% \arrayrulecolor{green!75!gray}
        \setlength{\unitlength}{\tabcolsep}
      \hspace{-5cm} \begin{tabular}[t]{|r|c|c|l|}
          \multicolumn{4}{l}{\C{main()}}\\ \hline
          ligne & a & b & Aff. \\ \hline
          initialisation  & -4.2 & 1.3 & \\ \hline
          6\pause & & &  \multicolumn{1}{ >{\columncolor{LightSteelBlue}}r|
          }{\put(2,0){\C{valeur\_absolue(-4.2)}}} \\ \hline
          \multicolumn{4}{|r|}{\only<7>{\cellcolor{LightGreen}}
            \put(1,1){\noindent
\rowcolors{2}{LightGoldenrod}{LightGoldenrod}
              \begin{tabular}[t]{|r|c|}\hline
                ligne  & x    \\ \hline
                init.  & -4.2 \\ \hline
                17 \pause & \multicolumn{1}{ >{\columncolor{LightSteelBlue}}r|
                }{\put(2,0){\C{maximum(-4.2, 4.2)}}}\\ \hline
                \multicolumn{2}{r|}{\only<7>{\cellcolor{LightGoldenrod}}
                  \put(1,1){\noindent
\rowcolors{2}{LightSkyBlue}{LightSkyBlue}
                    \begin{tabular}[t]{|r|c|c|}\hline
                      ligne & x & y \\ \hline
                      init.  & -4.2 & 4.2 \\ \hline
                      14\pause &\multicolumn{2}{|>{\columncolor{MediumSpringGreen}}l|}{renvoie 4.2}\\ \hline
                    \end{tabular}
                  }} \\ \hline
                17 \pause &\multicolumn{1}{|>{\columncolor{MediumSpringGreen}}l|}{renvoie 4.2}\\ \hline
              \end{tabular}
            }} \\ \hline
          \rowcolor{LightSkyBlue} 6 & & &\multicolumn{1}{ >{\columncolor{LightSkyBlue}}r| }{\put(2,0){\C{valeur\_absolue(1.3)}}}\\ \hline
          \multicolumn{4}{|r|}{\only<7>{\cellcolor{LightGreen}}
            \put(1,1){\noindent
\rowcolors{2}{LightGoldenrod}{LightGoldenrod}
              \begin{tabular}[t]{|r|c|}\hline
                ligne  & x    \\ \hline
                init.  & 1.3 \\ \hline
                17 & \multicolumn{1}{ >{\columncolor{LightSteelBlue}}r|
                }{\put(2,0){\C{maximum(1.3, -1.3)}}}\\ \hline
                \multicolumn{2}{|r|}{\only<7>{\cellcolor{LightGoldenrod}}
                  \put(1,1){\noindent
\rowcolors{2}{LightSkyBlue}{LightSkyBlue}
                    \begin{tabular}[t]{|r|c|c|}\hline
                      ligne & x & y \\ \hline
                      init.  & 1.3 & -1.3 \\ \hline
                      12 &\multicolumn{2}{|>{\columncolor{MediumSpringGreen}}l|}{renvoie 1.3}\\ \hline
                    \end{tabular}
                  }} \\ \hline
                17 &\multicolumn{1}{|>{\columncolor{MediumSpringGreen}}l|}{renvoie 1.3}\\ \hline
              \end{tabular}
            }} \\ \hline
          6  & 5.5 &  &  \\ \hline
          7  &  &  & \C{5.5} \\ \hline
          8 &\multicolumn{3}{|l|}{SUCCÈS}\\ \hline
        \end{tabular}
}%fin only
  \end{column}
\end{columns}
\end{frame}
\subsection{Traces : la mémoire et le temps}
\begin{frame}
  \frametitle{Traces : la mémoire  et le temps}
La trace d'un programme donne schématiquement ce type de dessin :
\begin{columns}
\column[T]{6.5cm}
\begin{itemize}
  \item<3-> verticalement, c'est le temps
  \item<3-> et horizontalement, l'occupation mémoire
  \item<4-> un appel de fonction occupe une portion de mémoire, puis la
    libère.
\item<5-> la trace représente réellement ce qui arrive dans vos
  programmes (durée de vie et localisation en mémoire des variables, etc.).
\end{itemize}
 \uncover<6->{
\structure{Un appel de fonction peut-il modifier la mémoire d'une
  fonction appelante ?}}
\column[T]{4.5cm}
\small
\vspace{-0.5cm}
\only<2->{\hspace{-0.7cm}\noindent\begin{tikzpicture}
\tikzstyle{every node}=[anchor=base]
\only<3->{\tikzset{axis/.style={fill=none}}}
\only<2>{\tikzset{axis/.style={color=white,fill=white,draw=white}}}
% Axis
  \draw[axis,arrows = ->] (-.7,5.9) -- (4.2,5.9);
  \draw(4.1,6) node[axis,anchor=base east] {mémoire};
  \draw[axis,arrows = ->] (-.6,6.1) -- (-.6,-.5);
  \draw(-.8,-.2) node[axis,rotate=-90,anchor=base east] {temps};
%
 \draw (0, 5.6) node {main()};
  \filldraw[fill=LightGreen,draw=Green] (-0.5,-0.25) rectangle
  (1,5.5);
  \draw[pattern color=Green, pattern=north east lines,draw=none] (-0.5,-0.25) rectangle
  (0,5.5);
 \draw (0.25, 5.25) node {\C{x}};
  \draw[draw=Green] (0.5,-0.25) -- (0.5,5.5);
  \draw (0.75, 5.25) node {\C{y}};
%
 \draw (1.75, 5.1) node {max(1,2)};
  \filldraw[fill=LightGoldenrod, draw=DarkGoldenrod] (1,4) rectangle (2.5,5);
  \draw[pattern color=DarkGoldenrod, pattern=north east lines,draw=none] (1,4) rectangle
  (1.5,5);
  \draw (2.25, 4.75) node {\C{b}};
  \draw[draw=DarkGoldenrod] (2,4) -- (2,5);
  \draw (1.75, 4.75) node {\C{a}};
%
  \draw (1.75, 3.6) node {\C{f(8)}};
  \filldraw[fill=PowderBlue,draw=RoyalBlue]
  (1,0) rectangle (2.5,3.5);
  \draw[pattern color=RoyalBlue, pattern=north east lines,draw=none]
  (1,0) rectangle  (1.5,3.5);
  \draw (1.75, 3.25) node {\C{z}};
  \draw[draw=RoyalBlue] (2,0) -- (2,3.5);
  \draw (2.25, 3.25) node {\C{r}};
%
  \draw (3, 3.1) node {\C{g(2)}};
  \filldraw[fill=Salmon,draw=Crimson] (2.5,2) rectangle(3.5,3);
  \draw[pattern color=Crimson, pattern=north east lines,draw=none]
  (2.5,2) rectangle  (3,3);
  \draw (3.25, 2.75) node {\C{n}};
%
  \filldraw[fill=LightGoldenrod, draw=DarkGoldenrod] (2.5,.5) rectangle
  (4,1.5);
  \draw[pattern color=DarkGoldenrod, pattern=north east lines,draw=none] (2.5,.5) rectangle
  (3,1.5);
  \draw (3.25, 1.6) node {max(2,-1)};
  \draw (3.25, 1.25) node {\C{a}};
  \draw[draw=DarkGoldenrod] (3.5,.5) -- (3.5,1.5);
  \draw (3.75, 1.25) node {\C{b}};
\end{tikzpicture}}
  \end{columns}
\end{frame}

\subsection{Pile d'appel}
\begin{frame}
  \frametitle{Pile d'appel}
\begin{columns}
\column[T]{4cm}
On parle de \alert{pile d'appel} car les appels de fonctions
s'empilent... \emph{comme sur une pile
    d'assiettes.}

\uncover<13->{
\structure{Peut-on avoir deux assiettes identiques dans la pile au
  même instant ? (La même
  fonction appelée deux fois \emph{simultanément})}}

% de contenus différents
\column[T]{7.5cm}
\small
\noindent\begin{tikzpicture}[rotate=90]
\tikzstyle{every node}=[anchor=base]
\tikzset{axis/.style={fill=none}}
\tikzset{fonction/.style={rotate=90}}
% Axis
  \draw[axis,arrows = ->] (-.7,5.9) -- (4.2,5.9);
  \draw(4.1,6) node[axis,rotate=90,anchor=base east] {mémoire};
  \draw[axis,arrows = ->] (-.6,6.1) -- (-.6,-.5);
  \draw(-.8,-.2) node[axis,anchor=base east] {temps};
%
\uncover<1,3->{
  \draw (0, 5.6) node[fonction] {main()};
  \filldraw[fill=LightGreen,draw=Green] (-0.5,-0.25) rectangle
  (1,5.5);
  \draw[pattern color=Green, pattern=north east lines,draw=none] (-0.5,-0.25) rectangle
  (0,5.5);
 \draw (0.25, 5.25) node {\C{x}};
  \draw[draw=Green] (0.5,-0.25) -- (0.5,5.5);
  \draw (0.75, 5.25) node {\C{y}};
}
%
\uncover<1,4,12->{
 \draw (1.75, 5.1) node[fonction] {max(1,2)};
  \filldraw[fill=LightGoldenrod, draw=DarkGoldenrod] (1,4) rectangle (2.5,5);
  \draw[pattern color=DarkGoldenrod, pattern=north east lines,draw=none] (1,4) rectangle
  (1.5,5);
  \draw (2.25, 4.75) node {\C{b}};
  \draw[draw=DarkGoldenrod] (2,4) -- (2,5);
  \draw (1.75, 4.75) node {\C{a}};
}
%
\uncover<1,6-10,12->{
  \draw (1.75, 3.6) node[fonction] {\C{f(8)}};
  \filldraw[fill=PowderBlue,draw=RoyalBlue]
  (1,0) rectangle (2.5,3.5);
  \draw[pattern color=RoyalBlue, pattern=north east lines,draw=none]
  (1,0) rectangle  (1.5,3.5);
  \draw (1.75, 3.25) node {\C{z}};
  \draw[draw=RoyalBlue] (2,0) -- (2,3.5);
  \draw (2.25, 3.25) node {\C{r}};
}
%
\uncover<1,7,12->{
  \draw (3, 3.1) node[fonction] {\C{g(2)}};
  \filldraw[fill=Salmon,draw=Crimson] (2.5,2) rectangle(3.5,3);
  \draw[pattern color=Crimson, pattern=north east lines,draw=none]
  (2.5,2) rectangle  (3,3);
  \draw (3.25, 2.75) node {\C{n}};
}
%
\uncover<1,9,12->{
  \filldraw[fill=LightGoldenrod, draw=DarkGoldenrod] (2.5,.5) rectangle
  (4,1.5);
  \draw[pattern color=DarkGoldenrod, pattern=north east lines,draw=none] (2.5,.5) rectangle
  (3,1.5);
  \draw (3.25, 1.6) node[fonction] {max(2,-1)};
  \draw (3.25, 1.25) node {\C{a}};
  \draw[draw=DarkGoldenrod] (3.5,.5) -- (3.5,1.5);
  \draw (3.75, 1.25) node {\C{b}};
}
\end{tikzpicture}
\end{columns}
\end{frame}

\begin{frame}[fragile]
  \frametitle{Factorielle récursive}
  \begin{block}{Définition}
    Une fonction récursive est une fonction dont la définition fait
    appel à la fonction \emph{elle-même}.
  \end{block}

\pause
Il y a une forte analogie avec les maths : $(n + 1)! = (n + 1) \times n!$

\pause
\begin{lstlisting}[escapechar={\%},basicstyle=\ttfamily\small]
int factorielle(int n)
{
    if (n < 2) { /* cas de base */
         return 1;
    }
    return n * factorielle(n - 1);
}
\end{lstlisting}
\end{frame}

\end{document}

% revenir sur :
%-l'appel au cours suivant : paramètre et expression, pile d'appel.
%-la grammaire des identificateurs
